\documentclass[11pt,a4paper]{article}
\usepackage{gtu-el-yazisi}
\begin{document}
\notbilgi{Analog Haberleşme}{ELM 361 Uygulama 3 -- SSB ve FM soruları}{ELM-361-Uygulama-3.pdf}

\sayfa{1}
\section{ELM 361 Uygulama -- 3}

\begin{ornek}[3.16 (SSB)]
An SSB-AM signal is generated by modulating an \qty{800}{\kilo\hertz} carrier by the signal $m(t)=\cos(2000\pi t)+2\sin(2000\pi t)$. The amplitude of the carrier is $A_c=100$.
\begin{enumerate}
\item Determine the signal $\hat m(t)$.
\item Determine the (time-domain) expression for the lower sideband of the SSB-AM signal.
\item Determine the magnitude spectrum of the lower-sideband SSB signal.
\end{enumerate}
\end{ornek}

\noindent\textbf{1.} $\hat m(t)=m_h(t)$ \quad hilbert transform.
\[ \hat m(t)=\cos\Big(2000\pi t-\tfrac{\pi}{2}\Big)+2\sin\Big(2000\pi t-\tfrac{\pi}{2}\Big) \]
\[ \hat m(t)=\sin(2000\pi t)-2\cos(2000\pi t) \]

\noindent\textbf{2.}
\[ y_{\mathrm{LSB}}(t)=A_c\big(m(t)\cos(2\pi f_c t)+\hat m(t)\sin(2\pi f_c t)\big) \]
\[ A_c=100,\quad f_c=800\,\mathrm{kHz},\quad m(t)=\cos(2000\pi t)+2\sin(2000\pi t) \]
\begin{align*}
y_{\mathrm{LSB}}(t)&=100\cdot\Big(\big(\cos(2000\pi t)+2\sin(2000\pi t)\big)\cdot\cos(2\pi\,800\mathrm{k}\cdot t)\\
&\qquad+\big(\sin(2000\pi t)-2\cos(2000\pi t)\big)\cdot\sin(2\pi\,800\mathrm{k}\cdot t)\Big)
\end{align*}
\begin{align*}
&=100\Big(\cos\big(\overbrace{2000\pi t}^{2\pi f_m}\big)\cos\big(2\pi\overbrace{800\mathrm{k}\,t}^{f_c}\big)+2\sin\big(\overbrace{2000\pi t}^{2\pi f_m}\big)\cos\big(2\pi\overbrace{800\mathrm{k}\,t}^{f_c}\big)\\
&\qquad+\sin\big(\overbrace{2000\pi t}^{2\pi f_m}\big)\sin\big(2\pi\overbrace{800\mathrm{k}\,t}^{f_c}\big)-2\cos\big(\overbrace{2000\pi t}^{2\pi f_m}\big)\sin\big(2\pi\overbrace{800\mathrm{k}\,t}^{f_c}\big)\Big)
\end{align*}
\formulkutu{\begin{aligned}f_c-f_m&=f_1\\ f_c+f_m&=f_2\end{aligned}}
\begin{align*}
&=100\Big(\tfrac{1}{2}\big(\cos(2\pi(f_m+f_c)t)+\cos(2\pi(f_m-f_c)t)\big)\;\text{\textcircled{\scriptsize 1}}\\
&\quad+\tfrac{2}{2}\big(\sin(2\pi(f_m+f_c)t)+\sin(2\pi(f_m-f_c)t)\big)\;\text{\textcircled{\scriptsize 2}}\,\text{\textcircled{\scriptsize 3}}\\
&\quad+\tfrac{1}{2}\big(\cos(2\pi(f_m-f_c)t)-\cos(2\pi(f_m+f_c)t)\big)\;\text{\textcircled{\scriptsize 1}}\\
&\quad-\tfrac{2}{2}\big(\sin(2\pi(f_c+f_m)t)+\sin(2\pi(f_c-f_m)t)\big)\Big)\;\text{\textcircled{\scriptsize 2}}\,\text{\textcircled{\scriptsize 3}}\\
&\qquad-\sin\big(2\pi(f_m-f_c)t\big)
\end{align*}

\sayfa{2}
\[ y_{\mathrm{LSB}}(t)=100\cdot\Big(\cos\big(2\pi(f_m-f_c)t\big)+2\sin\big(2\pi(f_m-f_c)t\big)\Big) \]
\[ =100\cdot\Big(\cos\big(2\pi(1000-800\mathrm{k})t\big)+2\sin\big(2\pi(1000-800\mathrm{k})t\big)\Big) \]
\[ y_{\mathrm{LSB}}(t)=100\Big(\cos\big(2\pi(-1000+\underbrace{800\mathrm{k}}_{f_c})t\big)-2\sin\big(2\pi(\underbrace{800\mathrm{k}}_{f_c}-1000)t\big)\Big) \]

\noindent\textbf{3.} $Y_{\mathrm{LSB}}(f)=\;?$

\noindent\textbf{Hint:} cos
\formulkutu{\begin{aligned}
\cos\alpha\cdot\cos\beta&=\tfrac12\big(\cos(\alpha+\beta)+\cos(\alpha-\beta)\big)\\
\cos\alpha\cdot\sin\beta&=\tfrac12\big(\sin(\beta+\alpha)+\sin(\beta-\alpha)\big)\\
\sin\alpha\cdot\cos\beta&=\tfrac12\big(\sin(\alpha-\beta)+\sin(\alpha+\beta)\big)\\
\sin\alpha\cdot\sin\beta&=\tfrac12\big(\cos(\alpha-\beta)-\cos(\alpha+\beta)\big)
\end{aligned}}

\begin{align*}
y_{\mathrm{LSB}}(t)&=100\cdot\Big(\cos\big(2\pi\underbrace{(f_c-f_m)}\,t\big)\\
&\quad-2\sin\big(2\pi(f_c-f_m)t\big)\Big)
\end{align*}
\begin{align*}
Y_{\mathrm{LSB}}(f)&=100\cdot\tfrac12\cdot\Big(\delta\big(f-(f_c-f_m)\big)+\delta\big(f+(f_c-f_m)\big)\Big)\\
&\quad-\frac{100}{{\not 2}\,j}\,2\Big(\delta\big(f-(f_c-f_m)\big)-\delta\big(f+(f_c-f_m)\big)\Big)
\end{align*}
\begin{align*}
Y_{\mathrm{LSB}}(f)&=50\cdot\Big(\delta(f-800\mathrm{k}+1000)+\delta(f+800\mathrm{k}-1000)\Big)\\
&\quad+100j\cdot\Big(\delta(f-800\mathrm{k}+1000)-\delta(f+800\mathrm{k}-1000)\Big)\\
&=(50+100j)\Big(\delta(f-800\mathrm{k}+1000)\Big)+(50-100j)\cdot\Big(\delta(f+800\mathrm{k}-1000)\Big)
\end{align*}
\noindent$\downarrow$ \textit{magnitude spectrum}
\begin{align*}
|Y_{\mathrm{LSB}}(f)|&=\sqrt{50^2+100^2}\cdot\Big(\delta(f-800\mathrm{k}+1000)+\delta(f+800\mathrm{k}-1000)\Big)\\
&=10\sqrt{125}\cdot\Big(\delta(f-800\mathrm{k}+1000)+\delta(f+800\mathrm{k}-1000)\Big)
\end{align*}

\sayfa{3}
\begin{ornek}[3.17]
Weaver's SSB modulator is illustrated in figure P-3.17. By taking the input signal as $m(t)=\cos 2\pi f_m t$ where \underline{$f_m<W$} demonstrate that by proper choice of $\kutula{f_1\ \text{and}\ f_2}$ the output is an SSB signal.
\end{ornek}
\elcizimi[0.9]{sekiller/uygulama-notlari-ELM-361-Uygulama-3-s3-1.png}[3]

\noindent\textit{input of upper LPF}
\begin{align*}
z_u(t)&=m(t)\cdot\cos(2\pi f_1 t)\\
&=\cos(2\pi f_m t)\cdot\cos(2\pi f_1 t)=\frac12\Big[\cos\big(2\pi\underbrace{(f_1-f_m)}_{f_a}t\big)+\cos\big(2\pi\underbrace{(f_1+f_m)}_{f_b}t\big)\Big]
\end{align*}
\noindent\textit{input of lower LPF}
\begin{align*}
z_\ell(t)&=\cos(2\pi f_m t)\cdot\sin(2\pi f_1 t)\\
&=\frac12\Big[\sin\big(2\pi(f_1-f_m)t\big)+\sin\big(2\pi(f_1+f_m)t\big)\Big]
\end{align*}

\noindent $\kutula{f_m<W}$ \quad $\longrightarrow$ \quad we will select $\kutula{f_1}$. \okunamadi{}
\formulkutu{|f_1-f_m|<\kutula{W}\quad\text{and}\quad f_1+f_m>W}
\okunamadi{W, W/2 (kutunun altındaki küçük notlar)} \hfill $\rightarrow$ SSB olması için seçilir.
\[ Z_u(f)=\frac14\cdot\Big(\delta(f-f_a)+\delta(f+f_a)\Big)+\frac14\cdot\Big(\delta(f-f_b)+\delta(f+f_b)\Big) \]
\elcizimi[0.5]{sekiller/uygulama-notlari-ELM-361-Uygulama-3-s3-2.png}[3]
\[ y_u(t)=\tfrac12\cos\big(2\pi(f_1-f_m)t\big) \]
\[ y_\ell(t)=\tfrac12\sin\big(2\pi(f_1-f_m)t\big) \]
\begin{align*}
y_{\mathrm{SSB}}(t)&=y_u(t)\cdot\cos(2\pi f_2 t)-y_\ell(t)\cdot\sin(2\pi f_2 t)\\
&=\tfrac12\cos\big(2\pi(f_1-f_m)t\big)\cdot\cos(2\pi f_2 t)-\tfrac12\sin\big(2\pi(f_1-f_m)t\big)\cdot\sin(2\pi f_2 t)
\end{align*}

\sayfa{4}
\begin{align*}
y_{\mathrm{SSB}}(t)&=\frac14\cdot\Big(\cos\big(2\pi(f_1-f_m+f_2)t\big)+\cos\big(2\pi(f_1-f_m-f_2)t\big)\Big)\\
&\quad-\frac14\cdot\Big(\cos\big(2\pi(f_1-f_m-f_2)t\big)-\cos\big(2\pi(f_1-f_m+f_2)t\big)\Big)
\end{align*}
\[ \cos(A-B)=\cos A\cos B+\sin A\sin B \]
\[ y_{\mathrm{SSB}}(t)=\tfrac12\cos\big(2\pi(f_1-f_m+f_2)t\big) \]
\elcizimi[0.4]{sekiller/uygulama-notlari-ELM-361-Uygulama-3-s4-1.png}[4]
\begin{align*}
f_1+f_2-f_m&=f_c+f_m\ \longrightarrow\ \overset{\text{at }f_c}{\text{USSB (upper)}}\\
f_1+f_2-f_m&=f_c-f_m\ \longrightarrow\ \text{LSSB (lower)}
\end{align*}
\noindent we are choosing different $f_1$ and $f_2$ for USSB and LSSB.

\noindent $f_c=\kutula{-(f_1+f_2)}$ \qquad $f_c=\kutula{f_1+f_2}$
\elcizimi[0.6]{sekiller/uygulama-notlari-ELM-361-Uygulama-3-s4-2.png}[4]

\sayfa{5}
\begin{ornek}[3.24 (FDM)]
We wish to transmit 60 voice-band signals by SSB (upper-sideband) modulation and freq. division multiplexing (FDM). Each of the 60 signals has a spectrum as shown figure P-3.24. Note that the voice-band signal is band limited to \qty{3}{\kilo\hertz}. If each signal is frequency translated separately, we require a frequency synthesizer that produces 60 carrier frequencies to perform the FDM. On the other hand, if we subdivide the channels into $L$ groups of $K$ subchannels each, such that $LK=60$; we may reduce the number of frequencies, from the synthesizer to $L+K$.
\begin{enumerate}[label=\arabic*)]
\item Illustrate the spectrum of the SSB signals in a group of $K$ subchannels. Assume that a \qty{1}{\kilo\hertz} guardband separates the signals in adjacent frequency subchannels and that the carrier frequencies are $f_{c1}=10\,\mathrm{kHz}$, $f_{c2}=14\,\mathrm{kHz}$ -- etc.
\item Sketch $L$ and $K$ such that $LK=60$ and $L+K$ is a minimum.
\item Determine the frequencies of the carriers if the 60 FDM signals occupy the frequency band $300$--$540\,\mathrm{kHz}$, and each group of $K$ signals occupies the band $10\,\mathrm{kHz}$ to $(10+4K)\,\mathrm{kHz}$.
\end{enumerate}
\end{ornek}
\elcizimi[0.8]{sekiller/uygulama-notlari-ELM-361-Uygulama-3-s5-1.png}[5]

\sayfa{6}
\noindent\textbf{1)}
\elcizimi[0.9]{sekiller/uygulama-notlari-ELM-361-Uygulama-3-s6-1.png}[6]

\noindent\textbf{2)} $\overset{\text{groups}}{L}\,\underset{\text{subchannels}}{K}=60$ \quad $(L,K)$ or $(K,L)$ $\rightarrow$ $(1,60),\ (2,30),\ (3,20),\ (4,15),\ (6,10)$.
\[ \text{Minimum value}\ \rightarrow\ \underset{\downarrow}{L}+\underset{\downarrow}{K} \]
\[ 6+10\ \rightarrow\ \underline{16} \]

\noindent\textbf{3)} $L=6,\ \underline{K=10}$ \quad we need 16 carriers with frequencies.

\noindent
\begin{minipage}[t]{0.52\linewidth}
\[ K=10\ \left\{\begin{array}{l}
f_{c1}=\kutula{10\,\mathrm{kHz}}\\
f_{c2}=14\,\mathrm{kHz}\\
f_{c3}=18\,\mathrm{kHz}\\
f_{c4}=22\,\mathrm{kHz}\\
\ \ \vdots\\
f_{c10}=46\,\mathrm{kHz}
\end{array}\right. \]
\[ \rlap{\rule[.55ex]{1.5em}{.4pt}}f_{c11}=\kutula{50\,\mathrm{kHz}}\quad 10+4K \]
\end{minipage}\hfill
\begin{minipage}[t]{0.44\linewidth}
\noindent $36\,\mathrm{kHz}$ \okunamadi{$+4\,\mathrm{kHz}$}
\[ L=\underline{6} \]
\fbox{$\begin{array}{l}
f_{L1}=\kutula{280\,\mathrm{kHz}}\\
f_{L2}=\kutula{320}\,\mathrm{kHz}\\
\qquad\ \ 360\\
\ \ \vdots\qquad\ 400\\
\qquad\ \ 440\\
f_{L6}=\kutula{480}\,\mathrm{kHz}\\
\rlap{\rule[.55ex]{8em}{.4pt}}f_{L7}=520\,\mathrm{kHz}
\end{array}$}
\[ \kutula{300}-540\,\mathrm{kHz} \]
\end{minipage}

\medskip
\noindent $L\rightarrow$ groups\\ $K\rightarrow$ subchannels.
\elcizimi[0.55]{sekiller/uygulama-notlari-ELM-361-Uygulama-3-s6-2.png}[6]
\[ \begin{array}{rl}
300 & \kutula{290}\\
340 & 330\\
380 & 370\\
420 & \vdots\\
460 & \\
500 & 490
\end{array} \]
\elcizimi[0.9]{sekiller/uygulama-notlari-ELM-361-Uygulama-3-s6-3.png}[6]
\noindent Eğer $f_{L1}$ grubunu 290'dan başlatırsak, subchannel'lar 300'den başlıyor olur.

\sayfa{7}
\begin{ornek}[4.3 (FM)]
The carrier $c(t)=100\cos(2\pi f_c t)$ is \underline{frequency modulated} by the signal $m(t)=5\cos(20000\pi t)$ where $f_c=10^8\,\mathrm{Hz}$. The peak frequency deviation is \underline{$20\,\mathrm{kHz}$}.
\begin{enumerate}[label=\arabic*.)]
\item Determine \underline{the amplitude and frequency of all signal components} that have a power-level of \underline{at least 10\%} of the power of the unmodulated carrier component.
\item from Carson's rule, determine the approximate BW of FM signal.
\end{enumerate}
\end{ornek}

\[ y(t)=A_c\sum_{n=-\infty}^{\infty}J_n(\beta)\cdot\cos\big(2\pi(f_c+n\cdot f_m)t\big) \]
\[ \begin{array}{ll} A_c=100 & f_m=10000\,\mathrm{Hz}\\ f_c=10^8\,\mathrm{Hz} & {}=10^4\,\mathrm{Hz} \end{array} \]
\begin{align*}
\beta&=\frac{\Delta f}{\underset{\downarrow}{B}}=\frac{k_f\cdot\overset{\max|m(t)|}{m_p}}{f_m}=\frac{\Delta f_{\max}}{f_m}\\
&\qquad\ \ f_m\ \ =\frac{20\cdot10^3}{10^4}=2
\end{align*}
\[ \beta=2. \]
\[ y(t)=\underline{100}\sum_{n=-\infty}^{\infty}\kutula{J_n(2)}\cdot\cos\big(2\pi(10^8+n\cdot10^4)t\big) \]

\noindent\textbf{1)} unmodulated carrier component $\rightarrow c(t)$
\[ P_c=\frac{A_c^2}{2}=\frac{100^2}{2}=5000 \]
\[ \text{at least }10\%\ \rightarrow\ 5000\cdot\frac{10}{100}=\kutula{500}\ \rightarrow\ \kutula{P_{J_n(2)}>500}\qquad P_{J_n(2)}=\frac{\big(A_c\,J_n(2)\big)^2}{2} \]
\begin{center}
\begin{tabular}{c|c|c|c|c}
Index $n$ & $J_n(2)$ & freq(Hz) & Amplitude $(100\,J_n(2))$ & Power $\rightarrow$\\ \hline
$0$ & .2239 & $10^8$ & 22.39 & 250,63\\
$\kutula{1}$ & .5767 & $10^8+10^4$ & 57.67 & \fbox{1663.1}\\
$\kutula{2}$ & .3528 & $10^8+2\cdot10^4$ & 35.28 & \fbox{622.46}\\
$3$ & .1289 & $10^8+3\cdot10^4$ & 12.89 & 83.13\\
$4$ & .0340 & $10^8+4\cdot10^4$ & 3.40 & 5.7785\\
\end{tabular}
\end{center}
\noindent Index $\rightarrow$ \okunamadi{$n=1,2$}\\
It's given $\rightarrow\times$ \qquad\quad $\times100\ \uparrow$

\sayfa{8}
\begin{center}
\begin{tabular}{c|c|c}
index & freq (Hz) & Amplitude\\ \hline
$1$ & $10^8+10^4$ & 57.67\\
$2$ & $10^8+2\cdot10^4$ & 35.28\\
$-1$ & $10^8-10^4$ & 57.67\\
$-2$ & $10^8-2\cdot10^4$ & 35.28
\end{tabular}
\end{center}
\[ J_n^2(2)=J_{-n}^2(2) \]
\noindent\textbf{2)}\quad $B_c=2\,\underset{\substack{\downarrow\\ f_m}}{B}\,(\beta+1)=2\cdot10^4(2+1)=\underline{6\cdot10^4\,\mathrm{Hz}}$

\sayfa{9}
\begin{ornek}[4.18 (FM)]
The modulation signal that is the input into an FM modulator is $m(t)=10\cos(16\pi t)$. The output of the FM modulator is
\[ y(t)=\underset{A_c}{10}\cos\Big[\overbrace{\underset{\substack{\downarrow\\ f_c}}{4000\pi t}+2\pi k_f\int_{-\infty}^{t}m(\tau)\,d\tau}^{\theta(t)}\Big] \]
where $k_f=10$ (See figure P-4.18). If the output of the FM modulator is passed through an ideal BPF centered at $f_c=2000$ with a bandwidth of $62\,\mathrm{Hz}$, determine the power of the frequency components at the output of the filter. What percentage of the transmitter power appears at the output of the BPF?
\end{ornek}
\elcizimi[0.9]{sekiller/uygulama-notlari-ELM-361-Uygulama-3-s9-1.png}[9]

\[ y(t)=A_c\sum_{n=-\infty}^{\infty}J_n(\beta)\cdot\cos\big(2\pi(f_c+n\cdot f_m)t\big) \]
\[ A_c=10,\quad f_c=2000,\quad f_m=8,\quad k_f=10 \]
\[ \underset{B\rightarrow}{\beta}=\frac{\Delta f}{f_m}=\frac{k_f\cdot m_p}{f_m}=\frac{k_f\cdot\max|m(t)|}{f_m}=\frac{10\cdot10}{8}=12.5 \]
\[ y(t)=10\cdot\sum_{n=-\infty}^{\infty}J_n(12.5)\cdot\cos\big(2\pi(2000+n\cdot8)t\big)\quad\text{at the output of the BPF.} \]
\[ \begin{array}{ll} & \text{freq}\\ J_0 & -\ 2000\\ J_1 & -\ 2008\\ J_2 & -\ 2016\\ J_3 & -\ 2024\\ J_4 & -\ 2032 \end{array}\qquad \begin{array}{ll} & \text{freq}\\ J_{-1} & -\ 1992\\ J_{-2} & -\ 1984\\ J_{-3} & -\ 1976\\ J_{-4} & -\ 1968 \end{array} \]
\[ \cdots\ -4\ \cdots\ \overset{n}{0}\ \cdots\ 4\ \cdots \]

\noindent$\rightarrow$ Power at the output of the BPF
\[ P_y=\frac{10^2}{2}\cdot J_0^2(12.5)+2\sum_{n=1}^{4}J_n^2(12.5)\cdot\frac{10^2}{2}=50\cdot0.2630 \]
\[ =\underline{13.15} \]
\begin{center}
\begin{tabular}{c|cccc}
$\beta$ & $J_0(\beta)$ & $J_1(\beta)$ & $J_2(\beta)$ & $\cdots$\\ \hline
1 & & & &\\
2 & & & &\\
3 & & & &\\
$\vdots$ & \textbf{Hint} & & &
\end{tabular}
\end{center}

\sayfa{10}
\[ \underset{\substack{\uparrow\\ \text{total}\\ \text{transmitted}\\ \text{power}}}{P_{tot}}=\frac{A_c^2}{2}=\frac{10^2}{2}=50. \]
\[ \frac{P_y}{P_{tot}}=\frac{13.15}{50}=\frac{26.30}{100}\rightarrow\underline{\%26.30} \]
\noindent Power at the output of BPF is only \%26.30 of transmitted power.

\noindent$\bullet$ Effective BW $\rightarrow$ Carson Rule \quad \%99\\
\%26.30 $\rightarrow$ \%99 olurdu

\sayfa{11}
\begin{ornek}[5.4.2]
A periodic square wave $m(t)$ (Fig. P.5.4-2a) frequency-modulates a carrier of frequency $f_c=\underline{10\,\mathrm{kHz}}$ with $\underline{\Delta f=1\,\mathrm{kHz}}$. The carrier amplitude is $A$. The resulting FM signal is demodulated, as shown in Fig. P.5.4-2.b. by the method discussed in sec. 5.4 (Fig. 5.11) sketch wave forms. b, c, d, e.
\end{ornek}
\elcizimi[0.9]{sekiller/uygulama-notlari-ELM-361-Uygulama-3-s11-1.png}[11]
\elcizimi[0.9]{sekiller/uygulama-notlari-ELM-361-Uygulama-3-s11-2.png}[11]
\begin{align*}
\Delta f&=k_f\cdot\overset{\max|m(t)|}{m_p}=k_f\\
&=1\mathrm{k}
\end{align*}
\begin{align*}
f_i&=f_c+k_f\cdot m(t)\\
f_i&=10\mathrm{k}+1\mathrm{k}\cdot m(t)\\
f_{i,\min}&=10\mathrm{k}+1\mathrm{k}\,\big(m(t)_{\min}\big)\\
f_{i,\min}&=9\,\mathrm{kHz}\\
f_{i,\max}&=10\mathrm{k}+1\mathrm{k}\,\big(m(t)_{\max}\big)\\
f_{i,\max}&=\underline{11\,\mathrm{kHz}}
\end{align*}
\elcizimi[0.7]{sekiller/uygulama-notlari-ELM-361-Uygulama-3-s11-3.png}[11]

\sayfa{12}
\noindent $b(t)$: \quad $y(t)=A_c\cdot\cos\Big(\underbrace{2\pi f_c t+2\pi k_f\int m(\tau)\,d\tau}_{\theta(t)}\Big)$

\noindent\textcircled{c}
\[ \frac{d}{dt}y(t)=A_c\cdot\theta'(t)\,\sin\big(\theta(t)\big) \]
\[ A_c\big[2\pi f_c+2\pi k_f\,m(t)\big]\cdot\sin\theta(t) \]
\elcizimi[0.5]{sekiller/uygulama-notlari-ELM-361-Uygulama-3-s12-1.png}[12]
\elcizimi[0.7]{sekiller/uygulama-notlari-ELM-361-Uygulama-3-s12-2.png}[12]
\elcizimi[0.9]{sekiller/uygulama-notlari-ELM-361-Uygulama-3-s12-3.png}[12]
\elcizimi[0.6]{sekiller/uygulama-notlari-ELM-361-Uygulama-3-s12-4.png}[12]

\sayfa{13}
\begin{ornek}[3.21]
A vestigial sideband modulation system is shown in figure P-3.21. The bandwidth of the message signal $m(t)$ is $W$, and the transfer func. of the band pass filter is shown in the figure.
\begin{enumerate}
\item Determine $h_\ell(t)$, which is the low pass equivalent of $h(t)$, where represents the impulse response of the band pass filter.
\item Derive an expression for the modulated signal $\underline{y(t)}$.
\end{enumerate}
\end{ornek}
\elcizimi[0.9]{sekiller/uygulama-notlari-ELM-361-Uygulama-3-s13-1.png}[13]
\[ M(f)=E(f)\cdot H_o(f) \]
\[ M(f)=M(f)\cdot\big(H_i(f-f_c)+H_i(f+f_c)\big)\cdot H_o(f) \]
\[ H_o(f)=\frac{1}{H_i(f-f_c)+H_i(f+f_c)},\qquad f\le B=W \]

\sayfa{14}
\elcizimi[0.9]{sekiller/uygulama-notlari-ELM-361-Uygulama-3-s14-1.png}[14]
\elcizimi[0.75]{sekiller/uygulama-notlari-ELM-361-Uygulama-3-s14-2.png}[14]
\[ M(f)=M(f)\cdot\big(H_i(f-f_c)+H_i(f+f_c)\big)\cdot H_o(f) \]
\elcizimi[0.35]{sekiller/uygulama-notlari-ELM-361-Uygulama-3-s14-3.png}[14]
\elcizimi[0.35]{sekiller/uygulama-notlari-ELM-361-Uygulama-3-s14-4.png}[14]

\end{document}
